Differentiating the moving integral gives $f(t+1)-f(t)=0$, so $1$ is a period. For example $f(x)=\sin2\pi x$ works and is not zero. For $g$, the second integral gives $g(a)=g(b)$. If both endpoint values vanish, there are already two zeros in $[a,b]$. Otherwise suppose they are positive (the negative case is analogous). Zero integral forces $g$ to be negative somewhere; continuity then gives one zero on each side of that point. Thus (i) is true. But (ii) is false: on $[0,1]$, $g(x)=\sin2\pi x$ has integral zero and equal endpoint values, yet only one interior zero, $\dfrac{1}{2}$.
